%%Rotat_rep_vent_moy
Application d'une rotation destin�e � annuler
les composantes transverses moyennes du vent :
\[
  \left(
\begin{array}{c}
u_r \\ v_r \\ w_r \\
\end{array}
  \right)
=
  \left(
\begin{array}{ccc}
\cos\epsilon \cos\theta & \sin\epsilon \cos\theta & \sin\theta \\
-\sin\epsilon & \cos\epsilon & 0 \\
-\cos\epsilon \sin\theta & -\sin\epsilon \sin\theta & \cos\theta \\
\end{array}
  \right)
.
  \left(
\begin{array}{c}
u \\ v \\ w \\
\end{array}
  \right)
\]

avec $\cos\epsilon=\frac{\overline{u}}{\sqrt{\overline{u}^2+\overline{v}^2}}$,
    $\sin\epsilon=\frac{\overline{v}}{\sqrt{\overline{u}^2+\overline{v}^2}}$,
    $\cos\theta=\frac{\sqrt{\overline{u}^2+\overline{v}^2}}
                     {\sqrt{\overline{u}^2+\overline{v}^2+\overline{w}^2}}$,
et  $\sin\theta=\frac{\overline{w}}
                     {\sqrt{\overline{u}^2+\overline{v}^2}+\overline{w}^2}$ 
(avec $\theta$ limit� � 20 degr�s).

Puis d�termination d'un angle $\beta$. minimisant $\overline{v'w'}$,
par r�solution it�rative de $2K\sin^2\beta+...= 0$

